\section{Experimental Setup}
\label{sec:setup}

\subsection{Implementation}
All models are implemented in PyTorch~\cite{paszke2019pytorch} and scikit-learn~\cite{pedregosa2011sklearn} and run on CPU with one thread per process. The data parser, grouping, causal features, agents, baselines, perturbation generator, protocol runner and statistics are released with the paper; one script regenerates every number from the raw XML files.

\subsection{Search Spaces and Budgets}
The search spaces (Supplementary Sec.~S2.1) are shared wherever a hyper-parameter is common to several families. Neural models: feature set (four options), trunk (64--64, 128--128, 256--256 or 256--512--256), dropout (0--0.3), normalization (batch or layer), learning rate ($3\times10^{-4}$--$3\times10^{-3}$), weight decay (0--$10^{-3}$) and batch size (256 or 512). Agents additionally: $\varepsilon$ decay per epoch (0.90--0.97), warm-up (0 or 4,000 random-action transitions) and updates per learning step (1, 2 or 4); E-DDQN: $\gamma\in\{0,0.5,0.9,0.99\}$, $\tau\in\{0.005,0.01,0.05\}$, replay composition $\rho\in\{0.4,0.5,0.6\}$ and full-information replay (on/off); E-CB: $\rho$ and full-information replay; DQN: $\gamma$ and the hard-update period (10, 100 or 500 updates); CS-MLP: uniform or class-balanced minibatches. All agents are trained for 40 epochs (one epoch is one pass over every training session in temporal order), with chunks of 256 transitions, a learning step every two chunks, five consolidation updates per epoch, a 50,000-transition buffer and gradient clipping at 1.0; the CS-MLP is trained for 40 epochs. Validation-based checkpoints are taken every two epochs. These budgets were calibrated on the October-2019 development-only sessions, where validation cost plateaued after 20--30 epochs. Each family receives \HPOBudget{} random configurations (12 for logistic regression), drawn once with a fixed seed and evaluated on every inner split of every outer fold. The same number of draws covers very different fractions of the search spaces: E-DDQN's space contains \PHSpaceEDDQN{} distinct configurations, the DQN's \PHSpaceDQN, E-CB's \PHSpaceECB{} and the CS-MLP's \PHSpaceCSMLP, against \PHSpaceHGB{} for gradient boosting, \PHSpaceRF{} for the random forest and \PHSpaceLR{} for logistic regression. The budget is therefore matched in the number of configurations, not in coverage; that a large RL search space is hard to explore with a small budget is itself a practical property of these detectors. Gradient boosting and random forests search learning rate, number of iterations, leaf counts, minimum leaf sizes, $\ell_2$ regularization, depth and feature sub-sampling; logistic regression searches the inverse regularization strength $C\in[10^{-4},10]$. The configuration selected for each outer fold is listed in Supplementary Sec.~S2.2.

\subsection{Seeds, Seed Means and the Seed Ensemble}
\label{sec:seeds}
Final models are trained with five seeds per outer fold; for every seed the outer-test score is the average of the \nInner{} cross-fitted inner models. We report mean~$\pm$~SD over seeds. Paired tests, bootstrap confidence intervals, the comparison with always alarming, event-level metrics, matched operating points, perturbation results and cost curves need a single prediction per sample; for these we use the \emph{seed ensemble}, i.e.\ the mean score over the five seeds, thresholded at the value chosen on the mean inner-OOF score. This choice was fixed in the evaluation code before the evaluation was run. Table~\ref{tab:main} reports both. They can differ, because thresholding an averaged score is not the same as averaging the metrics of separately thresholded seeds: for E-DDQN the ensemble has a higher 5:1 cost than the average seed (\PHEnsEDDQNCost{} vs.\ \EDDQNCostCostMean), whereas for the standard DQN it is lower (\PHEnsDQNCost{} vs.\ \DQNCostCostMean). Cost differences in the paired tests are therefore differences between seed ensembles and can be larger than the differences of seed means (E-DDQN vs.\ DQN: \PHEnsDiffDQN{} vs.\ \PHMeanDiffDQN). No seed was selected or discarded.

\subsection{Split-Protocol Comparison}
To quantify how the split protocol affects measured performance, we evaluated the same 21 raw columns of the curated release under three protocols with an untuned gradient-boosting classifier: (i) oversampling every discomfort grade to the size of grade~0 (5,243 samples) \emph{before} a random 70/15/15 row-level split, the protocol of the earlier, unpublished version of this study (it reproduces that version's partition sizes, 14,680 / 3,146 / 3,146, exactly); (ii) a row-level split without oversampling; and (iii) the grouped split of Section~\ref{sec:groups}. The earlier E-DDQN implementation was also run, unmodified, on partitions (i) and (iii). Because that version's 41-feature input files were not available, this is a reconstruction of its protocol rather than an exact reproduction.
